\documentclass[12pt]{article}
\usepackage{amsmath,amssymb,amsfonts}
\usepackage[margin=1in]{geometry}

\begin{document}

\textbf{Chapter 8, Problem 26.} Show that if $K$ is defined by Eq.\ (8.93) and $P_N$ by Eq.\ (8.94), then $K_N = P_NKP_N$.

\bigskip
\textbf{Solution.}

From the chapter, the operator $K$ is a completely continuous operator on a Hilbert space $H$ with orthonormal basis $\{\phi_i\}$. The projection operator $P_N$ projects onto the first $N$ basis elements:
\[
P_N f = \sum_{i=1}^{N}\langle\phi_i,f\rangle\,\phi_i,
\]
and the finite-rank approximation $K_N$ to $K$ is defined as:
\[
K_N = P_N K P_N.
\]

We verify this explicitly. Writing $K$ in terms of its matrix elements $K_{ij} = \langle\phi_i, K\phi_j\rangle$:
\[
Kf = \sum_{i,j}\phi_i\,K_{ij}\,\langle\phi_j,f\rangle.
\]

Now compute $P_NKP_N$:
\[
P_Nf = \sum_{j=1}^{N}\langle\phi_j,f\rangle\,\phi_j.
\]
Applying $K$:
\[
KP_Nf = K\sum_{j=1}^{N}\langle\phi_j,f\rangle\,\phi_j = \sum_{j=1}^{N}\langle\phi_j,f\rangle\,K\phi_j = \sum_{j=1}^N\langle\phi_j,f\rangle\sum_{i=1}^\infty K_{ij}\phi_i.
\]
Applying $P_N$:
\[
P_NKP_Nf = \sum_{j=1}^N\langle\phi_j,f\rangle\sum_{i=1}^N K_{ij}\phi_i = \sum_{i,j=1}^{N}K_{ij}\langle\phi_j,f\rangle\,\phi_i.
\]

This is precisely $K_N f$---the operator that acts like $K$ within the $N$-dimensional subspace spanned by $\{\phi_1,\ldots,\phi_N\}$ and gives zero on the orthogonal complement. It is a finite-rank operator (rank at most $N$) that approximates $K$.

By Theorem 8.4 of the chapter, $K_N\to K$ in operator norm as $N\to\infty$, since $K$ is completely continuous. This confirms that:
\[
\boxed{K_N = P_NKP_N.} \quad\blacksquare
\]

\end{document}
